\subsection{\bilingual{Result on the asymptotic speed of the left-most particle}{关于最左侧粒子渐近速率的结论}}

\noindent {\bf\Cref{theorem:speed}.}
\bilingual{
    For $a\in[0,1]$, consider a branching random walk $\mathcal V^{N|a}$ that starts with $\lfloor N^a\rfloor$ particles at position $0$, and whose branching law $\Xi$ satisfies \Cref{assumption:psi-zero,assumption:at-least-one-child,assumption:psi-one,assumption:psi-prime-one,assumption:finite-variance,assumption:log-expt-cross-sum-exp-neg-sum,assumption:moment-1,assumption:moment-4}. Then the speed of the leftmost particle in $\mathcal V^{N|a}$ exists and is given by
}{
    对于 $a\in[0,1]$,~ 考虑从位置 $0$ 处的 $\lfloor N^a\rfloor$ 个粒子出发的分支随机游走 $\mathcal V^{N|a}$,~ 其分支律 $\Xi$ 满足\Cref{assumption:psi-zero,assumption:at-least-one-child,assumption:psi-one,assumption:psi-prime-one,assumption:finite-variance,assumption:log-expt-cross-sum-exp-neg-sum,assumption:moment-1,assumption:moment-4}. ~则 $\mathcal V^{N|a}$ 中最左侧粒子的速率存在且由下式给出:
}
{
    \setlength\belowdisplayskip{-10pt}
    \begin{equation*}
        v_N
        =
        \lim_{n\to\infty}
        \frac{1}{n}\mathcal M^{N|a}_n
        =
        \frac{\pi^2}2\sigma^2\frac{1}{\log^2 N}
        +
        o_N\left(\frac{1}{\log^2 N}\right).
    \end{equation*}
}

\begin{proof}
    \quad

    \bilingual{
        The existence of the speed is guaranteed by \Cref{corollary:leftmost-speed:N0} under all assumptions. We proceed to calculate the speed by showing that there exists an upper bound and a lower bound of the same order. For both bounds, we are to construct, with respect to $\mathcal V^{N|a}$, a branching random walk, $\mathcal V^{\mathrm{upper}}$ and $\mathcal V^{\mathrm{lower}}$ respectively, such that the minimum position in $\mathcal V^{\mathrm{upper}}$ is point-wise (or stochastically) larger than the minimum position in question, and the minimum position in $\mathcal V^{\mathrm{lower}}$ is point-wise (or stochastically) smaller than the minimum position in question. We then calculate the speed of the left-most particle in $\mathcal V^{\mathrm{upper}}$ and $\mathcal V^{\mathrm{lower}}$ using the proved result, \Cref{proposition:3.1}, to get the desired upper and lower bounds respectively.
    }{
        速率的存在性由\Cref{corollary:leftmost-speed:N0} 保证. ~现在通过构造上界和下界来计算速率. ~对于这两个界,~ 分别构造与 $\mathcal V^{N|a}$ 耦合的分支随机游走 $\mathcal V^{\mathrm{upper}}$ 和 $\mathcal V^{\mathrm{lower}}$,~ 使得 $\mathcal V^{\mathrm{upper}}$ 中的最小位置在点态大于（或随机占优于）所讨论的最小位置,~ 而 $\mathcal V^{\mathrm{lower}}$ 中的最小位置在点态小于（或随机受制于）所讨论的最小位置. ~然后使用已证明的结果\Cref{proposition:3.1} 来计算 $\mathcal V^{\mathrm{upper}}$ 和 $\mathcal V^{\mathrm{lower}}$ 中最左侧粒子的速率,~ 从而得到所需的上界和下界.
    }

    \bilingual{
        For any $t>0$, choose and fix some sequence $\ell_N=\ell_N(t)$ such that
    }{
        对于任意 $t>0$,~ 选择并固定某个序列 $\ell_N=\ell_N(t)$ 使得
    }
    \begin{equation*}
        \lim_{N\to\infty}\frac{\ell_N}{\log^3N}=t.
    \end{equation*}

    \bilingual{
        For the upper bound, we construct, with respect to $\mathcal V^{N|a}$, an upper branching random walk $\mathcal V^{\mathrm{upper}}$ as follows:
    }{
        对于上界,~ 构造与 $\mathcal V^{N|a}$ 耦合的上界分支随机游走 $\mathcal V^{\mathrm{upper}}$ 如下:
    }

    \begin{figure}[H]
        \centering
        \begin{tikzpicture}[
            xscale=2,
            yscale=1.2,
            >=Stealth,
            % 样式定义
            survivor/.style={circle, fill=blue!70!black, inner sep=1.2pt},
            killed-node/.style={circle, draw=blue!70!black, fill=white, inner sep=1.2pt, cross out, thick, minimum size=4pt},
            edge-live/.style={thick, blue!70!black},
            edge-killed/.style={thick, dashed, blue!30},
            cloud-tri/.style={fill=blue!10, opacity=0.5},
            ref-line/.style={thin, gray!40, dashed}
        ]

            % 1. 时间周期参考线 (0, l_N, 2l_N, 3l_N)
            \foreach \x/\label in {0/0, 2/\ell_N, 4/2\ell_N, 6/3\ell_N} {
                \draw[ref-line] (\x, -0.5) -- (\x, 4.5);
                \node[below] at (\x, -0.5) {$\label$};
            }

            % ================= 第一周期: 0 到 l_N =================
            % 起点
            \node[survivor] (n0) at (0, 0) {};
            \node[left] at (n0) {$\emptyset$};

            % 三角形阴影 (顶点在 n0)
            % 假设在 l_N 时刻，粒子分布在 [0.8, 2.2] 之间
            \fill[cloud-tri] (0,0) -- (2, 0.8) -- (2, 2.2) -- cycle;

            % 高频分支
            \draw[edge-live] (0,0) -- (0.5, 0.2) -- (1.0, 0.5) -- (1.5, 0.4) -- (2, 0.8);
            \draw[edge-killed] (0.5, 0.2) -- (1.2, 1.2) -- (2, 1.3);
            \draw[edge-killed] (1.0, 0.5) -- (1.6, 1.0) -- (2, 1.8);
            \draw[edge-killed] (0,0) -- (0.8, 0.9) -- (2, 2.2);

            % 瓶颈处粒子
            \node[survivor] (m1) at (2, 0.8) {};
            \node[killed-node] at (2, 1.3) {};
            \node[killed-node] at (2, 1.8) {};
            \node[killed-node] at (2, 2.2) {};

            % ================= 第二周期: l_N 到 2l_N =================
            % 起点是 m1 (2, 0.8)
            % 假设在 2l_N 时刻，粒子分布在 [1.8, 3.2] 之间
            \fill[cloud-tri] (2, 0.8) -- (4, 1.8) -- (4, 3.2) -- cycle;

            % 高频分支
            \draw[edge-live] (2, 0.8) -- (2.5, 1.1) -- (3.0, 1.4) -- (3.5, 1.3) -- (4, 1.8);
            \draw[edge-killed] (2.5, 1.1) -- (3.2, 2.0) -- (4, 2.3);
            \draw[edge-killed] (3.0, 1.4) -- (3.6, 2.1) -- (4, 2.8);
            \draw[edge-killed] (2, 0.8) -- (2.8, 2.2) -- (4, 3.2);

            % 瓶颈处粒子
            \node[survivor] (m2) at (4, 1.8) {};
            \node[killed-node] at (4, 2.3) {};
            \node[killed-node] at (4, 2.8) {};
            \node[killed-node] at (4, 3.2) {};

            % ================= 第三周期: 2l_N 到 3l_N =================
            % 起点是 m2 (4, 1.8)
            % 假设在 3l_N 时刻，粒子分布在 [2.7, 4.1] 之间
            \fill[cloud-tri] (4, 1.8) -- (6, 2.7) -- (6, 4.1) -- cycle;

            % 高频分支
            \draw[edge-live] (4, 1.8) -- (4.5, 2.0) -- (5.0, 2.4) -- (5.5, 2.3) -- (6, 2.7);
            \draw[edge-killed] (4.5, 2.0) -- (5.2, 2.9) -- (6, 3.2);
            \draw[edge-killed] (5.0, 2.4) -- (5.6, 3.1) -- (6, 3.7);
            \draw[edge-killed] (4, 1.8) -- (4.8, 3.2) -- (6, 4.1);

            % 瓶颈处粒子
            \node[survivor] (m3) at (6, 2.7) {};
            \node[killed-node] at (6, 3.2) {};
            \node[killed-node] at (6, 3.7) {};
            \node[killed-node] at (6, 4.1) {};

            % ================= 标注 =================
            \node[below, blue!70!black] at (2, 0.8) {$\mathcal{M}^N_{\ell_N}$};
            \node[below, blue!70!black] at (4, 1.8) {$\mathcal{M}^N_{2\ell_N}$};
            \node[below, blue!70!black] at (6, 2.7) {$\mathcal{M}^N_{3\ell_N}$};

        \end{tikzpicture}
        \caption{\bilingual{Illustration of the construction of the upper branching random walk \texorpdfstring{$\mathcal{V}^{\mathrm{upper}}$}{Vupper}}{上界分支随机游走 \texorpdfstring{$\mathcal{V}^{\mathrm{upper}}$}{Vupper} 的构造示意图}}
    \end{figure}

    \begin{itemize}
        \item \bilingual{
            At time $0$, we kill all particles but one in the original $N$-BRW $\mathcal V^{N|a}$, so that there is only one particle at time $0$ and position $0$ in $\mathcal V^{\mathrm{upper}}$;
        }{
            在时刻 $0$,~ 除了一个粒子外,~ 在最初的 $N $-分支随机游走 $\mathcal V^{N|a}$ 中灭杀所有粒子,~ 使得 $\mathcal V^{\mathrm{upper}}$ 在零时刻只有一个粒子且位置为零;
        }
        \item \bilingual{
            From time $0$ till $\ell_N$, we let the process evolve as an $N$-BRW, i.e., at each step, we keep the left-most $N$ particles and kill all other particles;
        }{
            从时刻 $0$ 到 $\ell_N$,~ 让过程以 $N$-分支随机游走的方式演化,~ 即在每一步,~ 保留最左侧的 $N$ 个粒子并灭杀所有其他粒子;
        }
        \item \bilingual{
            At time $\ell_N$, we kill all particles but one left-most particle in $\mathcal V^{\mathrm{upper}}$;
        }{
            在时刻 $\ell_N$,~ 除了最左侧的一个粒子外,~ 在 $\mathcal V^{\mathrm{upper}}$ 中灭杀所有粒子;
        }
        \item \bilingual{
            From time $\ell_N$ onward, we let the process evolve as an $N$-BRW and repeat the killing procedure at all times $k\ell_N$ for $k\in\mathbb N$.
        }{
            从时刻 $\ell_N$ 开始,~ 让过程以 $N$-分支随机游走的方式演化,~ 并在所有时刻 $k\ell_N$ 处重复上述灭杀程序,~ 其中 $k\in\mathbb N$.
        }
    \end{itemize}
    \bilingual{
        By definition of $\mathcal V^{\mathrm{upper}}$, we have, by \Cref{corollary:2} and by induction on $k$, on the same probability space,
    }{
        根据 $\mathcal V^{\mathrm{upper}}$ 的定义有,~ 由\Cref{corollary:2} 和对 $k$ 的归纳法,~ 在相同的概率空间上,
    }
    % \begin{equation*}
    $
        \mathcal V^{\mathrm{upper}}_{k\ell_N}
        \succcurlyeq
        \mathcal V^{N|a}_{k\ell_N},
        \,
        % \quad
        \forall k\in\mathbb N,
    % \end{equation*}
    $
    \bilingual{
        and, by some coupling of $\mathcal V^{\mathrm{upper}}$ and $\mathcal V^{N|a}$,
    }{
        并且,~ 通过 $\mathcal V^{\mathrm{upper}}$ 和 $\mathcal V^{N|a}$ 的某个耦合,
    }
    \begin{equation*}
        \mathcal V^{\mathrm{upper}}_n
        \succcurlyeq
        \mathcal V^{N|a}_n,
        \quad\forall n\in[\![k\ell_N,(k+1)\ell_N]\!],\,\forall k\in\mathbb N.
    \end{equation*}
    \bilingual{
        Hence, there is some coupling of $\mathcal V^{\mathrm{upper}}$ and $\mathcal V^{N|a}$ such that
    }{
        因此,~ 存在 $\mathcal V^{\mathrm{upper}}$ 和 $\mathcal V^{N|a}$ 的某个耦合使得
    }
    % \begin{equation*}
    $
        \mathcal M^{\mathrm{upper}}_n
        \ge
        \mathcal M^{N|a}_n,
        % \quad
        \,
        \forall n\in\mathbb N.
    % \end{equation*}
    $
    \bilingual{
        Furthermore, by definition of $\mathcal V^{\mathrm{upper}}$, we have
    }{
        此外,~ 根据 $\mathcal V^{\mathrm{upper}}$ 的定义有
    }
    \begin{equation*}
        \mathcal V^{\mathrm{upper}}_{(k+1)\ell_N+n}
        -
        \mathcal M^{\mathrm{upper}}_{k\ell_N}
        \stackrel{\mathrm{d}}{=}
        \mathcal V^{N|0}_{\ell_N},
        \quad k\in\mathbb N,
    \end{equation*}
    \bilingual{
        and they are independent for different $k$'s, where, both sides of the equality are regarded as random finite counting measures, and the right-hand side is the population measure obtained from an $N$-BRW $\mathcal V^{N|0}$ starting with $N^0=1$ particle at position $0$ on the same probability space. Hence,
    }{
        并且对于不同的 $k$,~ 它们是独立的,~ 其中等式两边都被视为随机有限计数测度,~ 右侧是从一个 $N$-分支随机游走 $\mathcal V^{N|0}$ 中得到的粒子分布,~ 该 $N$-分支随机游走在相同的概率空间上以 $N^0=1$ 个粒子在位置 $0$ 处开始. ~因此,
    }
    \begin{equation*}
        \left(
            \min\left(
                \mathcal V^{\mathrm{upper}}_{(k+1)\ell_N}
                -
                \mathcal M^{\mathrm{upper}}_{k\ell_N}
            \right):
            k\in\mathbb N
        \right)
    \end{equation*}
    \bilingual{
        is an i.i.d.~sequence of random variables with the same distribution as $\min\mathcal V^{N|0}_{\ell_N}$. To use the Strong Law of Large Numbers (SLLN), we need to check that, given $N$,
    }{
        是一个独立同分布的随机变量序列,~ 其分布与 $\min\mathcal V^{N|0}_{\ell_N}$ 的分布相同. ~为了使用大数定律,~ 需要检查对于给定的 $N$,
    }
    \begin{equation*}
        \mathbb E\left[
            \min\mathcal V^{N|0}_{\ell_N}
        \right]
        =
        \mathbb E\left[
            \mathcal M^{N|0}_{\ell_N}
        \right]
        <
        \infty.
    \end{equation*}
    \bilingual{
        Indeed, under \Cref{assumption:psi-one,assumption:moment-1}, i.e., $\psi(1)=0$ and $\mathbb E[(\Xi_1)_+]<\infty$, we have
    }{
        事实上,~ 在\Cref{assumption:psi-one,assumption:moment-1} 下,~ 即 $\psi(1)=0$ 和 $\mathbb E[(\Xi_1)_+]<\infty$,~ 有
    }
    \begin{equation*}
        \left|
            \mathbb E\left[
                    \mathcal M^{N|0}_n
            \right]
        \right|
        \le
        \mathbb E\left[
            \left(
                \mathcal M^{N|0}_n
            \right)_+
        \right]
        +
        \mathbb E\left[
            \left(
                \mathcal M^{N|0}_n
            \right)_-
        \right]
        <
        \infty,
    \end{equation*}
    \bilingual{
        see \Cref{lemma:M:moment-1:neg,lemma:M:moment-1:pos} under the assumptions for details. Hence, by SLLN and the proven result of \Cref{proposition:3.1} for $a=0$,
    }{
        见\Cref{lemma:M:moment-1:neg,lemma:M:moment-1:pos} 在假设下的详细说明. ~因此,~ 通过大数定律和\Cref{proposition:3.1} 已证明的 $a=0$ 的结果,
    }
    \begin{align*}
        v_N
        &=
        \lim_{n\to\infty}
        \frac{1}{n}\mathcal M^{N|a}_n
        =
        \lim_{k\to\infty}
        \frac 1{k\ell_N}
        \mathcal M^{N|a}_{k\ell_N}
        \\ &
        \le
        \limsup_{k\to\infty}
        \frac 1{k\ell_N}
        \mathcal M^{\mathrm{upper}}_{k\ell_N}
        % \\ &
        =
        \frac 1{\ell_N}
        \limsup_{k\to\infty}
        \frac 1{k}
        \left(
            \sum_{i=1}^k
            \left(
                \mathcal M^{\mathrm{upper}}_{i\ell_N}
                -
                \mathcal M^{\mathrm{upper}}_{(i-1)\ell_N}
            \right)
        \right)
        \\ &
        =
        \frac 1{\ell_N}
        \limsup_{k\to\infty}
        \frac 1{k}
        \left(
            \sum_{i=1}^k
            \min
            \left(
                \mathcal V^{\mathrm{upper}}_{i\ell_N}
                -
                \mathcal M^{\mathrm{upper}}_{(i-1)\ell_N}
            \right)
        \right)
        \\ &
        =
        \frac 1{\ell_N}
        \mathbb E\left[
            \min\mathcal V^{N|0}_{\ell_N}
        \right]
        % \\ &
        =
        \frac 1{\ell_N}
        \mathbb E\left[
            \mathcal M^{N|0}_{\ell_N}
        \right].
    \end{align*}

    \bilingual{
        For the last expectation in the above, by the proved result of \Cref{proposition:3.1} and under the additional \Cref{assumption:moment-exp-1}, it holds for $a=0$ that \eqref{eq:4.13} is true, i.e.,
    }{
        对于这里的最后一项中的期望,~ 根据上文证得的\Cref{proposition:3.1},~ 在补充的\Cref{assumption:moment-exp-1} 下,~ 对于$a=0$的情形成立 \eqref{eq:4.13},~ 即
    }
    \begin{equation*}
        \mathbb E\left[
            \mathcal M_{\ell_N}^{N|0}
        \right]
        =
        \frac{\pi^2}2\sigma^2t\log N
        +
        o_N(\log N),
    \end{equation*}
    \bilingual{
        so that we can get the desired upper bound by plugging it into the above. However, this assumption is not necessary, and we will directly prove the original formula without the condition of \Cref{assumption:moment-exp-1} by using the restarting mechanism used in the proof of the upper bound for $a=0$ (Section \ref{section:a:0}). To this end, since we can let $t\to\infty$ eventually, we can instead show that
    }{
        代入原式即可得证速率上界. ~但这一假设是不必要的,~ 下面将在去除该\Cref{assumption:moment-exp-1} 的条件下,~ 利用\Cref{proposition:3.1} 的$a=0$情形的上界证明 (第 \ref{section:a:0} 节) 中用到的重启机制,~ 直接证明原式成立. 为此,~ 由于最终可以令$t\to\infty$,~ 可以转而证明
    }
    \begin{equation*}
        \mathbb E\left[
            \mathcal M_{\ell_N}^{N|0}
        \right]
        \le
        \frac{\pi^2}2\sigma^2t\log N
        +
        O_N(\log N),
    \end{equation*}
    \bilingual{
        and thus
    }{
        从而有
    }
    \begin{align*}
        v_N
        &
        \le
        \frac 1{\ell_N}\mathbb E\left[
            \mathcal M_{\ell_N}^{N|0}
        \right]
        \\
        &
        \le
        \frac 1{t\log^3N+o_N(\log^3N)}
        \left[
            \frac{\pi^2}2\sigma^2t\log N
            +
            O_N(\log N)
        \right]
        \\[5pt] &
        \le
        \inf_{t>0}\left[
            \left(
                \frac{O_N(1)}t
                +
                \frac{\pi^2}2\sigma^2
            \right)
            \frac{1}{\log^2 N}
            +
            o_N\left(\frac{1}{\log^2 N}\right)
        \right]
        \\[5pt] &
        =
        \frac{\pi^2}2\sigma^2\frac{1}{\log^2 N}
        +
        o_N\left(\frac{1}{\log^2 N}\right).
    \end{align*}
    \bilingual{
        Indeed, we use the restarting mechanism specified in Section \ref{section:a:0} and use the same notation as in that section. Let $\mathcal V^{\mathrm b}$ be the branching random walk constructed in that section from time $0$ to time $\ell_N$, and let
    }{
        事实上,~ 使用第 \ref{section:a:0} 节中用到的重启机制,~ 并使用该节中的符号. ~令 $\mathcal V^{\mathrm b}$ 是在该节中从时刻 $0$ 到时刻 $\ell_N$ 构造的分支随机游走,~ 并令
    }
    \begin{equation*}
        E_N
        :=
        \left\{
            \kappa\le K,\,
            \tau\le T
        \right\},
    \end{equation*}
    \bilingual{
        where $\kappa$ is the number of restarts and $\tau$ is the total time of restarts, while $K=\lfloor\delta\log N\rfloor$ and $T=\gamma\log N$ for some $\delta>0$ and $\gamma>0$. We have
    }{
        其中 $\kappa$ 是重启的次数,~ $\tau$ 是重启的总时间,~ 而 $K=\lfloor\delta\log N\rfloor$ 和 $T=\gamma\log N$ 其中 $\delta>0$ 和 $\gamma>0$. ~有
    }
    \begin{align}
        &\quad
        \mathbb E\left[
            \mathcal M_{\ell_N}^{N|0}
        \right]
        % \\
        % &
        \le
        \mathbb E\left[
            \mathcal M_{\ell_N}^{\mathrm b}
        \right]
        \nonumber \\ &
        =
        \mathbb E\left[
            \left(
                \mathcal M_{\ell_N}^{\mathrm b}
                -
                \max\mathcal V^{\mathrm b}_\tau
            \right)
            \mathbf 1_{E_N}
        \right]
        +
        \mathbb E\left[
            \max\mathcal V^{\mathrm b}_\tau
            \mathbf 1_{E_N}
        \right]
        +
        \mathbb E\left[
            \mathcal M_{\ell_N}^{\mathrm b}
            \mathbf 1_{E_N^{\mathrm c}}
        \right]
        \label{eq:velocity:restart:decomp}
    \end{align}
    \bilingual{
        For the first term of the above equation, since $\tau$ is a stopping time, by \nameref{lemma:branching-property-stopping-time} and construction of $\mathcal V^{\mathrm b}$, we have
    }{
        对于上式的第一项,~ 由于 $\tau$ 是一个停时,~ 根据\nameref{lemma:branching-property-stopping-time}和 $\mathcal V^{\mathrm b}$ 的构造,~ 有
    }
    \begin{align*}
        \mathbb E\left[
            \left(
                \mathcal M_{\ell_N}^{\mathrm b}
                -
                \max\mathcal V^{\mathrm b}_\tau
            \right)
            \mathbf 1_{E_N}
        \right]
        &
        =
        \mathbb E\left[
            \mathbb E\left(
                \mathcal M_{\ell_N}^{\mathrm b}
                -
                \max\mathcal V^{\mathrm b}_\tau
                \middle|
                \mathcal F_\tau
            \right)
            \mathbf 1_{E_N}
        \right]
        \\ &
        =
        \mathbb E\left[
            \left.\mathbb E\left[
                \mathcal M_{\ell_N-\iota}^{N|\alpha}
            \right]\right|_{\iota=\tau}
            \mathbf 1_{E_N}
        \right]
        \\ &
        =
        \mathbb E\left[
            \left.\left(
                \left(
                    -\alpha
                    +
                    \frac{\pi^2}2\sigma^2t
                \right)\log N
                +
                o_N(\log N)
            \right)\right|_{\iota=\tau}
            \mathbf 1_{E_N}
        \right]
        \\ &
        \le
        \left(
            -\alpha
            +
            \frac{\pi^2}2\sigma^2t
        \right)\log N
        +
        o_N(\log N),
    \end{align*}
    \bilingual{
        where the third equality is by the proved result of \Cref{proposition:3.1} for $a=\alpha$. For the second term of \eqref{eq:velocity:restart:decomp}, we have, by \eqref{eq:a:0:max-Zb-tau},
    }{
        其中第三个等号是由\Cref{proposition:3.1} 已证明的 $a=\alpha$ 的结果得出. ~对于 \eqref{eq:velocity:restart:decomp} 的第二项,~ 根据 \eqref{eq:a:0:max-Zb-tau},~ 有
    }
    \begin{align*}
        \mathbb E\left[
            \max\mathcal V^{\mathrm b}_\tau
            \mathbf 1_{E_N}
        \right]
        &
        \le
        \mathbb E\left[
            \left(S_\kappa+(\kappa-1)M+\ell M\right)
            \mathbf 1_{E_N}
        \right]
        \\ &
        \le
        \mathbb E\left[
            S_\kappa
            \mathbf 1_{E_N}
        \right]
        +
        (K-1)M
        +
        \ell M
        \\ &
        =
        \mathbb E\left[S_\kappa\right]
        + O_N(\log N),
    \end{align*}
    \bilingual{
        where $S_\kappa$ is the sum of $\kappa$ i.i.d.~copies of $X_1$, $\kappa$ is a geometric random variable with parameter $\mathbb P(\#\mathcal Z^{(M)}\in[N^\alpha,N])\ge p_1>0$, $X_1$ is the sum of $G_1$ i.i.d.~copies of $\Xi_1$, and $G_1$ is a geometric random variable with parameter $p_0=\mathbb P(\#\Xi^{(M)\mathrm b}\ge 2)>0$. We can show by Wald's identity that
    }{
        其中 $S_\kappa$ 是 $\kappa$ 个独立同分布的 $X_1$ 的和,~ $\kappa$ 是一个参数为 $\mathbb P(\#\mathcal Z^{(M)}\in[N^\alpha,N])\ge p_1>0$ 的几何随机变量,~ $X_1$ 是 $G_1$ 个独立同分布的 $\Xi_1$ 的和,~ 而 $G_1$ 是一个参数为 $p_0=\mathbb P(\#\Xi^{(M)\mathrm b}\ge 2)>0$ 的几何随机变量. 通过 Wald 等式可以得到
    }
    \begin{align*}
        \mathbb E\left[S_\kappa\right]
        &
        =
        \mathbb E\left[\kappa\right]
        \mathbb E\left[X_1\right]
        =
        \mathbb E\left[\kappa\right]
        \mathbb E\left[G_1\right]
        \mathbb E\left[\Xi_1\right]
        \\ &
        \le
        \frac 1{p_1}
        \frac{1-p_0}{p_0}
        \left(
            \mathbb E\left[(\Xi_1)_+\right]
            +
            \mathbb E\left[(\Xi_1)_-\right]
        \right)
        =
        O_N(1),
    \end{align*}
    \bilingual{
        where the last equality is by \Cref{assumption:moment-1} (i.e. $\mathbb E[(\Xi_1)_+]<\infty$) and \Cref{lemma:moment-1:neg} under \Cref{assumption:psi-one} (i.e. $\psi(1)=0$). For the third and last term of \eqref{eq:velocity:restart:decomp}, we estimate it using the Cauchy--Schwarz inequality. By \Cref{lemma:M:moment-4:neg,lemma:M:moment-4:pos}, we have $\mathbb E[(\mathcal M_n^N)^2] = O(n^2)$ (since the second moment is controlled by the fourth moment). Given $\ell_N = O(\log^3 N)$, we have $\mathbb E[(\mathcal M_{\ell_N}^{\mathrm b})^2]^{1/2} = O(\log^3 N)$. We can combine it with the polynomial decay $\mathbb P(E_N^{\mathrm c}) \le N^{-\varepsilon}$ from \eqref{eq:kappa:bound} and \eqref{eq:kappa:tau:bound}. Indeed, under the requirement \eqref{eq:gamma:condition} on $\gamma$ that
    }{
        其中最后一个等式是由\Cref{assumption:moment-1} (即 $\mathbb E[(\Xi_1)_+]<\infty$) 和\Cref{lemma:moment-1:neg} (在\Cref{assumption:psi-one} 下, 即 $\psi(1)=0$) 得出. ~对于 \eqref{eq:velocity:restart:decomp} 的第三项也是最后一项,~ 利用 Cauchy--Schwarz 不等式进行估计. ~根据\Cref{lemma:M:moment-4:neg,lemma:M:moment-4:pos},~ 我们有 $\mathbb E[(\mathcal M_n^N)^2] = O(n^2)$ (二阶矩由四阶矩控制). ~考虑到 $\ell_N = O(\log^3 N)$,~ 成立 $\mathbb E[(\mathcal M_{\ell_N}^{\mathrm b})^2]^{1/2} = O(\log^3 N)$. ~结合 \eqref{eq:kappa:bound} 和 \eqref{eq:kappa:tau:bound} 中可推得的概率多项式衰减 $\mathbb P(E_N^{\mathrm c}) \le N^{-\varepsilon}$. ~事实上,~ 在$\gamma$满足 \eqref{eq:gamma:condition} 的要求时,~ 即
    }
    \begin{equation*}
        \gamma
        >
        \frac{\delta}{\theta}
        \mathbb E\left[
            e^{\theta G_1}
        \right]
        +
        \delta
        +
        \frac{\alpha}{\log m},
    \end{equation*}
    \bilingual{
        we have
    }{
        有
    }
    \begin{align*}
        \mathbb P\left(E_N^{\mathrm c}\right)
        &
        =
        \mathbb P\left(
            \kappa>K
        \right)
        +
        \mathbb P\left(
            \kappa\le K,\,\tau>T
        \right)
        % \\ &
        \le
        N^{-\varepsilon_1}
        +
        N^{-\varepsilon_2}
        \le
        N^{-\varepsilon},
    \end{align*}
    \bilingual{
        which holds for $\varepsilon:=\frac 12(\varepsilon_1\wedge\varepsilon_2)>0$ and $N$ big enough. Therefore,
    }{
        对 $\varepsilon:=\frac 12(\varepsilon_1\wedge\varepsilon_2)>0$ 和足够大的 $N$ 成立. ~因此,
    }
    \begin{align*}
        \left| \mathbb E\left[ \mathcal M_{\ell_N}^{\mathrm b} \mathbf 1_{E_N^{\mathrm c}} \right] \right|
        \le \mathbb E\left[ (\mathcal M_{\ell_N}^{\mathrm b})^2 \right]^{1/2} \mathbb P(E_N^{\mathrm c})^{1/2}
        \le O(\log^3 N) \cdot N^{-\varepsilon/2}
        = o_N(\log N).
    \end{align*}
    \bilingual{
        Finally, the upper bound for $v_N$ is proved by plugging these estimates back into \eqref{eq:velocity:restart:decomp} and letting $N \to \infty$ followed by $t \to \infty$.
    }{
        最后,~ 将上述估计代入 \eqref{eq:velocity:restart:decomp},~ 并令 $N \to \infty$ 继而 $t \to \infty$,~ 即可得证 $v_N$ 上界.
    }

    \bilingual{
        For the lower bound, we construct, by modifying $\mathcal V^{N|1}$ (i.e. an $N$-BRW starting with $N^1=N$ particles at position $0$ on the same probability space such that $\mathcal V^{N|a}$ is a subtree of $\mathcal V^{N|1}$), a lower branching random walk $\mathcal V^{\mathrm{lower}}$ for $\mathcal V^{N|a}$ as follows:
    }{
        对于下界,~ 通过修改 $\mathcal V^{N|1}$ (即在相同的概率空间上以 $N^1=N$ 个粒子在位置 $0$ 处开始的 $N$-分支随机游走,~ 使得 $\mathcal V^{N|a}$ 是 $\mathcal V^{N|1}$ 的一个子树),~ 构造 $\mathcal V^{N|a}$ 的下界分支随机游走 $\mathcal V^{\mathrm{lower}}$ 如下:
    }

    \begin{figure}[htbp]
        \centering
        \begin{tikzpicture}[
            xscale=2,
            yscale=1.2,
            >=Stealth,
            % 样式定义
            survivor/.style={circle, fill=blue!70!black, inner sep=1.2pt},
            other-node/.style={circle, draw=blue!70!black, fill=white, inner sep=1pt},
            edge-live/.style={thick, blue!70!black},
            edge-killed/.style={thick, dashed, blue!30},
            cloud-tri/.style={fill=blue!10, opacity=0.5},
            ref-line/.style={thin, gray!40, dashed},
            translation-arrow/.style={->, red!80!black, thin, dashed}
        ]

            % 1. 时间周期参考线 (0, l_N, 2l_N, 3l_N)
            \foreach \x/\label in {0/0, 2/\ell_N, 4/2\ell_N, 6/3\ell_N} {
                \draw[ref-line] (\x, -0.5) -- (\x, 4.5);
                \node[below] at (\x, -0.5) {$\label$};
            }

            % ================= 第一周期: 0 到 l_N =================
            % 起点 (N个粒子在0处，表现为一个点)
            \node[survivor] (n0) at (0, 0) {};
            \node[left] at (n0) {$\emptyset$};

            % 三角形阴影
            \fill[cloud-tri] (0,0) -- (2, 0.8) -- (2, 2.2) -- cycle;

            % 高频分支演化
            \draw[edge-live] (0,0) -- (0.4, 0.2) -- (0.8, 0.6) -- (1.4, 0.5) -- (2, 0.8);
            \draw[edge-live] (0,0) -- (0.6, 0.8) -- (1.2, 1.1) -- (2, 1.4);
            \draw[edge-live] (0,0) -- (0.5, 0.9) -- (1.5, 1.8) -- (2, 2.2);

            % 时刻 l_N 的粒子分布
            \node[survivor] (m1) at (2, 0.8) {}; % 最小值点
            \node[survivor] at (2, 1.4) {};
            \node[survivor] at (2, 2.2) {};

            % 平移过程示意 (将上方粒子压到最小值位置)
            \draw[translation-arrow] (2, 2.2) to[bend left] (2, 0.9);
            \draw[translation-arrow] (2, 1.4) to[bend left] (2, 0.85);

            % ================= 第二周期: l_N 到 2l_N =================
            % 全部 N 个粒子从 m1 (2, 0.8) 出发
            \fill[cloud-tri] (2, 0.8) -- (4, 1.8) -- (4, 3.2) -- cycle;

            % 高频分支
            \draw[edge-live] (2, 0.8) -- (2.5, 1.1) -- (3.0, 1.6) -- (3.5, 1.5) -- (4, 1.8);
            \draw[edge-live] (2, 0.8) -- (2.6, 1.8) -- (3.2, 2.2) -- (4, 2.5);
            \draw[edge-live] (2, 0.8) -- (2.7, 2.1) -- (3.6, 2.9) -- (4, 3.2);

            % 时刻 2l_N 粒子分布
            \node[survivor] (m2) at (4, 1.8) {}; % 最小值点
            \node[survivor] at (4, 2.5) {};
            \node[survivor] at (4, 3.2) {};

            % 平移示意
            \draw[translation-arrow] (4, 3.2) to[bend left] (4, 1.9);
            \draw[translation-arrow] (4, 2.5) to[bend left] (4, 1.85);

            % ================= 第三周期: 2l_N 到 3l_N =================
            % 从 m2 (4, 1.8) 出发
            \fill[cloud-tri] (4, 1.8) -- (6, 2.7) -- (6, 4.1) -- cycle;

            \draw[edge-live] (4, 1.8) -- (4.5, 2.1) -- (5.1, 2.5) -- (6, 2.7);
            \draw[edge-live] (4, 1.8) -- (4.8, 2.9) -- (5.5, 3.4) -- (6, 4.1);
            \draw[edge-live] (4, 1.8) -- (4.6, 2.5) -- (5.3, 3.0) -- (6, 3.4);

            % 时刻 3l_N 最小值点
            \node[survivor] (m3) at (6, 2.7) {};
            \node[survivor] at (6, 3.4) {};
            \node[survivor] at (6, 4.1) {};

            % ================= 标注 =================
            \node[below, blue!70!black] at (2, 0.8) {$\mathcal{M}^N_{\ell_N}$};
            \node[below, blue!70!black] at (4, 1.8) {$\mathcal{M}^N_{2\ell_N}$};
            \node[below, blue!70!black] at (6, 2.7) {$\mathcal{M}^N_{3\ell_N}$};

        \end{tikzpicture}
        \caption{\bilingual{Illustration of the construction of the lower branching random walk \texorpdfstring{$\mathcal{V}^{\mathrm{lower}}$}{Vlower}}{下界分支随机游走 \texorpdfstring{$\mathcal{V}^{\mathrm{lower}}$}{Vlower} 构造示意图}}
    \end{figure}

    \begin{itemize}
        \item \bilingual{
            At time $0$, we keep all particles in the original $N$-BRW $\mathcal V^{N|1}$ and let the process evolve as $\mathcal V^{N|1}$ from $0$ till $\ell_N$;
        }{
            在时刻 $0$,~ 保留最初的 $N$-分支随机游走 $\mathcal V^{N|1}$ 中的所有粒子,~ 让过程以 $\mathcal V^{N|1}$ 的方式从 $0$ 演化到 $\ell_N$;
        }
        \item \bilingual{
            At time $\ell_N$, even if the population size $\#\mathcal Z^{N|a}_{\ell_N}$ of $\mathcal V^{N|a}$ does not reach $N$, we translate all the $N$ particles in $\mathcal V^{\mathrm{lower}}_{\ell_N}$ to the left-most position $\min\mathcal V^{\mathrm{lower}}_{\ell_N}$;
        }{
            在时刻 $\ell_N$,~ 即使 $\mathcal V^{N|a}$ 的粒子总数 $\#\mathcal Z^{N|a}_{\ell_N}$ 没有达到 $N$,~ 也将 $\mathcal V^{\mathrm{lower}}_{\ell_N}$ 中的所有 $N$ 个粒子平移到最左侧位置 $\min\mathcal V^{\mathrm{lower}}_{\ell_N}$;
        }
        \item \bilingual{
            From time $\ell_N$ onward, we let the process evolve as $\mathcal V^{N|1}$ and repeat the translation procedure at all times $k\ell_N$ for $k\in\mathbb N$.
        }{
            从时刻 $\ell_N$ 开始,~ 让过程以 $\mathcal V^{N|1}$ 的方式演化,~ 并在所有时刻 $k\ell_N$ 对 $k\in\mathbb N$ 重复平移过程.
        }
    \end{itemize}
    \bilingual{
        By definition of $\mathcal V^{\mathrm{lower}}$, we have, by \Cref{corollary:2} and by induction on $k$, on the same probability space,
    }{
        根据 $\mathcal V^{\mathrm{lower}}$ 的定义有,~ 由\Cref{corollary:2} 和对 $k$ 的归纳法,~ 在相同的概率空间上,
    }
    % \begin{equation*}
    $
        \mathcal V^{N|a}_{k\ell_N}
        \succcurlyeq
        \mathcal V^{\mathrm{lower}}_{k\ell_N},
        \quad\forall k\in\mathbb N,
    % \end{equation*}
    $
    \bilingual{
        and, by some coupling of $\mathcal V^{\mathrm{lower}}$ and $\mathcal V^{N|a}$,
    }{
        并且,~ 通过 $\mathcal V^{\mathrm{lower}}$ 和 $\mathcal V^{N|a}$ 的某个耦合,
    }
    \begin{equation*}
        \mathcal V^{N|a}_n
        \succcurlyeq
        \mathcal V^{\mathrm{lower}}_n,
        \quad\forall n\in[\![k\ell_N,(k+1)\ell_N]\!],\,\forall k\in\mathbb N.
    \end{equation*}
    \bilingual{
        Hence, there is some coupling of $\mathcal V^{\mathrm{lower}}$ and $\mathcal V^{N|a}$ such that
    }{
        因此,~ 存在 $\mathcal V^{\mathrm{lower}}$ 和 $\mathcal V^{N|a}$ 的某个耦合,~ 使得
    }
    % \begin{equation*}
    $
        \mathcal M^{N|a}_n
        \ge
        \mathcal M^{\mathrm{lower}}_n,
        % \quad
        \,
        \forall n\in\mathbb N.
    % \end{equation*}
    $
    \bilingual{
        Furthermore, by definition of $\mathcal V^{\mathrm{lower}}$, we have
    }{
        此外,~ 根据 $\mathcal V^{\mathrm{lower}}$ 的定义有
    }
    \begin{equation*}
        \mathcal V^{\mathrm{lower}}_{(k+1)\ell_N}
        -
        \mathcal M^{\mathrm{lower}}_{k\ell_N}
        \stackrel{\mathrm{d}}{=}
        \mathcal V^{N|1}_{\ell_N},
        \quad k\in\mathbb N,
    \end{equation*}
    \bilingual{
        and they are independent for different $k$'s. Hence,
    }{
        并且对于不同的 $k$,~ 它们是独立的. ~因此,
    }
    \begin{equation*}
        \left(
            \min\left(
                \mathcal V^{\mathrm{lower}}_{(k+1)\ell_N}
                -
                \mathcal M^{\mathrm{lower}}_{k\ell_N}
            \right):
            k\in\mathbb N
        \right)
    \end{equation*}
    \bilingual{
        is an i.i.d.~sequence of random variables with the same distribution as $\min\mathcal V^{N|1}_{\ell_N}$. To use SLLN, we need to check that, given $N$,
    }{
        是一个独立同分布的随机变量序列,~ 其分布与 $\min\mathcal V^{N|1}_{\ell_N}$ 的分布相同. ~为了使用大数定律,~ 需要检查对于给定的 $N$,
    }
    \begin{equation*}
        \mathbb E\left[
            \min\mathcal V^{N|1}_{\ell_N}
        \right]
        =
        \mathbb E\left[
            \mathcal M^{N|1}_{\ell_N}
        \right]
        <
        \infty.
    \end{equation*}
    \bilingual{
        Indeed, since $\mathcal M^{N|1}_n\ge\mathcal M^{N|0}_n$ holds (stochastically) for all $n$, we have
    }{
        事实上,~ 由于 $\mathcal M^{N|1}_n\le\mathcal M^{N|0}_n$ 对所有 $n$ 都 (随机地) 成立,~ 有
    }
    \begin{equation*}
        \mathbb E\left[
            \left(
                \mathcal M^{N|1}_n
            \right)_+
        \right]
        \le
        \mathbb E\left[
            \left(
                \mathcal M^{N|0}_n
            \right)_+
        \right]
        <
        \infty,
    \end{equation*}
    \bilingual{
        and
    }{
        并且
    }
    \begin{align*}
        \mathbb E\left[
            \left(
                \mathcal M^{N|1}_n
            \right)_-
        \right]
        &
        \le
        \mathbb E\left[
            e^{
                \left(
                    \mathcal M^{N|1}_n
                \right)_-
            }
        \right]
        % \\ &
        \le
        \mathbb E\left[
            e^{
                -\mathcal M^{N|1}_n
            }
            +
            1
        \right]
        \\ &
        \le
        1
        +
        \sum_{|u|=n}
        \mathbb E\left[
            e^{
                -\mathcal V^{N|1}(u)
            }
        \right]
        % \\ &
        =
        1
        +
        N\mathbb E[W_n]
        =
        1+N<\infty,
    \end{align*}
    \bilingual{
        where $W_n$ is the additive martingale, and $\mathbb E[W_n]=1$ is proven in \Cref{lemma:M:exp-neg-moment} under \Cref{assumption:psi-one}, i.e., $\psi(1)=0$. Hence, by SLLN and the proven result of \Cref{proposition:3.1} for $a=1$,
    }{
        其中 $W_n$ 是加性鞅,~ $\mathbb E[W_n]=1$ 在\Cref{assumption:psi-one} 下由\Cref{lemma:M:exp-neg-moment} 证明,~ 即 $\psi(1)=0$. ~因此,~ 通过大数定律和\Cref{proposition:3.1} 已证明的 $a=1$ 的结果,
    }
    \begin{align*}
        v_N
        &=
        \lim_{n\to\infty}
        \frac{1}{n}\mathcal M^{N|a}_n
        =
        \lim_{k\to\infty}
        \frac 1{k\ell_N}
        \mathcal M^{N|a}_{k\ell_N}
        \\ &
        \ge
        \liminf_{k\to\infty}
        \frac 1{k\ell_N}
        \mathcal M^{\mathrm{lower}}_{k\ell_N}
        % \\ &
        =
        \frac 1{\ell_N}
        \liminf_{k\to\infty}
        \frac 1{k}
        \left(
            \sum_{i=1}^k
            \left(
                \mathcal M^{\mathrm{lower}}_{i\ell_N}
                -
                \mathcal M^{\mathrm{lower}}_{(i-1)\ell_N}
            \right)
        \right)
        \\ &
        =
        \frac 1{\ell_N}
        \liminf_{k\to\infty}
        \frac 1{k}
        \left(
            \sum_{i=1}^k
            \min
            \left(
                \mathcal V^{\mathrm{lower}}_{i\ell_N}
                -
                \mathcal M^{\mathrm{lower}}_{(i-1)\ell_N}
            \right)
        \right)
        \\ &
        =
        \frac 1{\ell_N}
        \mathbb E\left[
            \min\mathcal V^{N|1}_{\ell_N}
        \right]
        % \\ &
        =
        \frac 1{\ell_N}
        \mathbb E\left[
            \mathcal M^{N|1}_{\ell_N}
        \right]
        \\ &
        =
        \frac 1{t\log^3N+o(\log^3N)}
        \left(
            \left(
                -1+
                \frac{\pi^2}2\sigma^2t
            \right)\log N
            +
            o_N(\log N)
        \right)
        \\ &
        =
        \left(
            -\frac 1t+
            \frac{\pi^2}2\sigma^2
        \right)
        \frac{1}{\log^2 N}
        +
        o_N\left(\frac{1}{\log^2 N}\right),
    \end{align*}
    \bilingual{
        and let $t\to\infty$ to get the desired lower bound. We are now done with the proof of \Cref{theorem:speed}.
    }{
        并让 $t\to\infty$ 来得到所需的下界. ~至此,~ \Cref{theorem:speed} 得证.
    }
\end{proof}
